\section{Wannier functions in practice}\label{app:E.5}

The first-principles phonon-assisted absorption spectrum of silicon was converged to an energy resolution of \SI{14}{meV} with $40\times40\times40$ grids for both the electron wavevector $\vect{k}$ and the phonon wavevector $\vect{Q}$~\cite{noffsinger2012phonon}. That is $64{,}000^2 \approx 4\times10^{9}$ pairs $(\vect{k},\vect{Q})$, each with a vertex $g_{mn\nu}(\vect{k},\vect{Q})$ for every pair of bands and every branch. A direct DFPT calculation costs about one ground-state calculation per perturbation and wavevector, so $64{,}000$ wavevectors are already out of reach, and the matrix elements between all those states even more so. The same study computed the vertices directly on a $6^3$ grid and produced the rest by Wannier interpolation. In practice the work is in choosing and testing the Wannier functions and in keeping every interpolated quantity in one gauge.

\subsection{Computed quantities}

A Wannier calculation turns quantities known on a coarse uniform grid into real-space matrices, and those back into the same quantities anywhere in the zone. The bands are interpolated through \cref{eq:3.6-HWR,eq:3.6-HWk}, and the eigenvectors of $H^W(\vect{k})$ define $U(\vect{k})$ on the fine grid. For the electron--phonon vertex, the coarse-grid vertices are rotated with the coarse-grid $U$ matrices and the phonon displacement patterns are undone, which leaves a matrix in two lattice vectors, one for the electron hop and one for the displaced atom~\cite[Eqs.~(177)--(179)]{giustino2017electron}. Back on a fine grid,
\begin{equation}
  g_{mn\nu}(\vect{k},\vect{Q}) = \sum_{\vect{R}_e,\vect{R}_p}\ee^{\ii(\vect{k}\cdot\vect{R}_e + \vect{Q}\cdot\vect{R}_p)}
  \sum_{m'n'\kappa\alpha}U_{mm'}(\vect{k}+\vect{Q})\,g^W_{m'n'\kappa\alpha}(\vect{R}_e,\vect{R}_p)\,U^\dagger_{n'n}(\vect{k})\,
  \frac{e^\alpha_{\kappa\nu}(\vect{Q})}{\sqrt{2M_\kappa\omega_{\vect{Q}\nu}}} ,
  \label{eq:E.5-ginterp}
\end{equation}
in the normalization of \cref{eq:A.6-g}. The phonon frequencies and polarization vectors on the fine grid come from the interpolated force constants of \cref{app:E.4}. The same construction interpolates the velocity matrix elements of \cref{sec:3.2}, which the optical limit needs, and, with the restriction of \cref{sec:3.6}, the density vertex $\rho_{mn}(\vect{k},\vect{q})$.

\subsection{Projections and disentanglement}

The four bond-centered Wannier functions of the valence bands of silicon need no disentanglement, and they are the natural choice for the valence part of a polarization calculation (\cref{eq:E.4-berry}). A dark-matter rate needs the conduction bands as well, and with them the eight disentangled $sp^3$ functions.

Projection onto atomic $sp^3$ orbitals alone reproduces the bands of silicon well except where higher bands with some $sp^3$ character mix in; forcing the whole valence manifold into the frozen window removes those deviations~\cite[Fig.~5]{marzari2012maximally}. For an electron-recoil calculation the inner (frozen) window should contain every state that can be a final state at the energies of interest, plus a margin, since a conduction state outside the frozen window is reproduced only approximately.

The number of Wannier functions grows with the energy range that must be described. For phonon-assisted absorption up to the direct gap and beyond, Ref.~\cite{noffsinger2012phonon} extracted 26 Wannier functions from 34 bands, reproducing the band structure from \SI{10}{eV} below to \SI{30}{eV} above the Fermi level. The range is set by the intermediate states. A second-order amplitude sums over all bands $m$, and the Wannier manifold truncates that sum; the remainder above the manifold can be added by the projected Sternheimer equation \cref{eq:E.4-dpsiperp}.

Every set of Wannier functions passes four checks before use. The spread functional $\Omega_W$ must be converged and the individual spreads small compared with the cell. The centers must sit where the symmetry says, on the bonds or the atoms. For a time-reversal-symmetric crystal the maximally localized functions can be chosen real, and a large ratio of imaginary to real parts signals a poor minimum. Above all, the interpolated bands must agree with a direct calculation at points off the coarse grid, inside the frozen window, to the tolerance the application needs. In silicon the last check has a twist: the conduction minimum lies on the $\Delta$ line at about 0.85 of the way from $\Gamma$ to X, not at a high-symmetry point, so a coarse grid that does not resolve it misplaces the valley. Ref.~\cite{ponce2018towards} needed a $12^3$ coarse $\vect{k}$ grid for the conduction bands for this reason, while $6^3$ was enough for the phonon wavevectors.

\subsection{Interpolating vertices, and the long-range parts}

The decay of $g^W(\vect{R}_e,\vect{R}_p)$ decides whether \cref{eq:E.5-ginterp} converges. In the electron vector $\vect{R}_e$ it decays at least as fast as the Wannier functions themselves. In the atom vector $\vect{R}_p$ it decays as the screened potential of a displaced atom. In a nonpolar semiconductor or insulator that potential falls off at least as fast as a quadrupole, $R_p^{-3}$, so the vertices in reciprocal space stay finite as $\vect{Q}\to0$ and the interpolation converges~\cite[Sect.~VI.A.2]{giustino2017electron}. For diamond, coarse grids of $4^3$, $6^3$ and $8^3$ points interpolate the vertex of the highest optical branch along high-symmetry lines onto the direct DFPT values, with the $4^3$ grid already close~\cite[Fig.~6]{giustino2017electron}. For a new material, the same comparison, of $g$ from direct DFPT at a few $(\vect{k},\vect{Q})$ off the coarse grid with the interpolated values, is the only evidence that the interpolation converges.

In a polar crystal the displaced atom carries a dipole, its potential decays as $R_p^{-2}$, and the vertex of the longitudinal optical branch diverges as $1/Q$: the Fr\"ohlich coupling. No finite real-space sum reproduces a divergence, so the long-range part is split off analytically. With the book's normalization,
\begin{equation}
  g^{\mathrm{L}}_{mn\nu}(\vect{k},\vect{Q}) = \ii\,\frac{e^2}{\Omega}\sum_\kappa\frac{1}{\sqrt{2M_\kappa\omega_{\vect{Q}\nu}}}
  \sum_{\vect{G}\ne-\vect{Q}}\frac{\vect{e}_{\kappa\nu}(\vect{Q})\cdot Z^*_\kappa\cdot(\vect{Q}+\vect{G})}{(\vect{Q}+\vect{G})\cdot\varepsilon_\infty\cdot(\vect{Q}+\vect{G})}\,
  \langle\psi_{m,\vect{k}+\vect{Q}}|\ee^{\ii(\vect{Q}+\vect{G})\cdot(\vect{r}-\vect{\tau}_\kappa)}|\psi_{n\vect{k}}\rangle ,
  \label{eq:E.5-gL}
\end{equation}
with $\vect{\tau}_\kappa$ the position of atom $\kappa$ in the cell. It is the potential of the bound charge $-\nabla\cdot\vect{P}$ of the displacement wave, with $\vect{P} = e\,\vect{u}_\kappa\cdot Z^*_\kappa/\Omega$ per cell, screened by $\varepsilon_\infty$ and acting on an electron of charge $-e$, for the displacement per unit $B_{\vect{Q}\nu}$ of \cref{eq:A.6-u}; it is Ref.~\cite[Eq.~(180)]{giustino2017electron} in Heaviside--Lorentz units and the normalization of \cref{app:A.6}~\cite{verdi2015frohlich}. The procedure subtracts $g^{\mathrm{L}}$ from the coarse-grid vertices, interpolates the smooth short-range remainder $g^{\mathrm{S}}$, and adds $g^{\mathrm{L}}$ back on the fine grid. The overlap in \cref{eq:E.5-gL} is needed at the fine-grid points; for small $\vect{Q}+\vect{G}$ it is $[U(\vect{k}+\vect{Q})U^\dagger(\vect{k})]_{mn}$ in the Wannier gauge~\cite[Sect.~VI.A.3]{giustino2017electron}, and the sum over $\vect{G}$ is cut off by a Gaussian factor in practice. The inputs $Z^*_\kappa$ and $\varepsilon_\infty$ come from \cref{app:E.4}, with the caveat about the functional's $\varepsilon_\infty$ stated there.

Besides its dipole, a displaced atom carries a dynamical quadrupole, which is nonzero even in silicon, where every Born charge vanishes, and its potential gives vertices that are finite as $\vect{Q}\to 0$ but depend on the direction of $\vect{Q}$. Such a nonanalytic angular dependence is long-ranged in real space and is not removed by a finer coarse grid: densifying the grid confines the spurious oscillations of the interpolated potential to a smaller region around $\Gamma$ without reducing their amplitude~\cite{brunin2020electron}. In silicon, omitting the quadrupole term changes the computed electron mobility by about \SI{10}{\percent}, and in GaAs by about \SI{30}{\percent}~\cite{brunin2020electron}. Codes that support it subtract the quadrupolar term along with the dipole; where the code does not, the error must be estimated from the small-$Q$ behavior of directly computed vertices.

Every quantity that enters one amplitude must be interpolated with the same $U(\vect{k})$. The safest practice is to compute all coarse-grid quantities from one set of wavefunctions, rotate them with one set of $U$ matrices, and interpolate them with one $H^W(\vect{R})$. Quasiparticle corrections are interpolated on the same Wannier basis, and since they change the energy denominators but not the wavefunctions, velocity matrix elements must be rescaled by the ratio of quasiparticle to Kohn--Sham transition energies, \cref{eq:E.3-vrescale}~\cite{noffsinger2012phonon}.

\subsection{Density vertices at large momentum}

In the Wannier gauge the density vertex $\rho_{mn}(\vect{k},\vect{q})$ is the matrix $\langle w_{m'\vect{0}}|\ee^{\ii\vect{q}\cdot\vect{r}}|w_{n'\vect{R}}\rangle$, Fourier transformed in $\vect{R}$. For small $\vect{q}$ one may expand $\ee^{\ii\vect{q}\cdot\vect{r}} \approx 1 + \ii\vect{q}\cdot\vect{r}$, and the vertex follows from the overlaps and the position matrix elements already computed for the Wannier functions; this is the small-$\vect{q}$ form used for $g^{\mathrm{L}}$ above. The expansion is controlled by $|\vect{q}|$ times the spatial extent of a Wannier function, which is of the order of a bond length, \SI{2.35}{\angstrom} in silicon. At the momentum transfers of electron recoil, a keV and above, $|\vect{q}|\times\SI{2.35}{\angstrom} \approx 1.2$ at \SI{1}{keV} and grows linearly beyond, so the expansion fails. The momentum transfer is a parameter of the vertex, not a lattice wavevector, and for each $\vect{q}$ of interest the real-space matrix must be built anew.

The practical route therefore computes $\rho_{mn}(\vect{k},\vect{q})$ directly on the coarse $\vect{k}$ grid for each target $\vect{q}$, from the plane-wave coefficients of the states at $\vect{k}$ and at $\vect{k}+\vect{q}$, rotates it with the same $U$ matrices into the Wannier gauge, transforms it to real space in $\vect{R}$, and interpolates in $\vect{k}$ alone. The cost is one non-self-consistent calculation at the shifted grid $\vect{k}+\vect{q}$ per target $\vect{q}$, and the number of targets is set by the resolution needed in $|\vect{q}|$ and in direction. For a cubic crystal the directions can be restricted to the irreducible wedge.

\subsection{Counting matrix elements and memory}

\begin{workedexample}{Matrix elements of a converged silicon rate}\label{we:E-counts}
The setup of Ref.~\cite{noffsinger2012phonon} has $N_W = 26$ Wannier functions, coarse $6^3$ grids for $\vect{k}$ and $\vect{Q}$, and fine $40^3$ grids for both, with two atoms and six branches. We count what a phonon-assisted rate needs and where it must live.

\emph{Fine-grid vertices.} There are $N_k N_Q = 64{,}000^2 = 4.1\times10^9$ pairs $(\vect{k},\vect{Q})$, and at each pair $N_W^2\times6 = 4{,}056$ complex vertices. At \SI{16}{bytes} per complex number this is $1.7\times10^{13}$ numbers, or \SI{270}{TB}. Nobody stores it. Codes loop over $\vect{Q}$ (or over $\vect{k}$), interpolate the vertices for all partners at once, accumulate the rate, and discard them. Even one $\vect{Q}$ with all $64{,}000$ partners and all band pairs is \SI{4}{GB}, which is why the band indices are restricted at once to the bands that can be initial, final or intermediate states at the energies of interest.

\emph{Real-space vertices.} The matrix $g^W(\vect{R}_e,\vect{R}_p)$ of \cref{eq:E.5-ginterp} has one entry for each pair of lattice vectors in the Wigner--Seitz supercells of the two coarse grids, about $6^3 = 216$ vectors each, times $N_W^2$ band pairs and six displacement directions: $216^2\times676\times6 \approx 1.9\times10^8$ complex numbers, \SI{3}{GB}. This is the object that is kept, and it must fit in the memory of one node, or be distributed over several.

\emph{Interpolation work.} \Cref{eq:E.5-ginterp} is done in two stages. For each $\vect{Q}$ the sum over $\vect{R}_p$ costs $216\times216\times676\times6 \approx 1.9\times10^8$ complex multiply--adds; for each $\vect{k}$ at that $\vect{Q}$ the sum over $\vect{R}_e$ costs $216\times676\times6 \approx 8.8\times10^5$, plus the $U$ rotations, which are smaller. Over all pairs the second stage dominates: $4.1\times10^9\times8.8\times10^5 \approx 3.6\times10^{15}$ multiply--adds, about $3\times10^{16}$ floating-point operations. Restricting $\vect{Q}$ to the irreducible wedge divides this by the 48 operations of the cubic group, to $6\times10^{14}$, and a pair of energy windows does far more.

\emph{Energy windows.} Near the absorption edge only valence states within a few tenths of an eV of the valence maximum and conduction states within the same range of the six conduction minima contribute. For a parabolic band of mass $m^*$, the states within $\delta$ of the edge fill a sphere of radius $k_\delta = (2m^*\delta)^{1/2}$, a fraction $(4\pi/3)k_\delta^3/V_{\mathrm{BZ}}$ of the zone, with $V_{\mathrm{BZ}} = (2\pi)^3/\Omega = \SI{6.19}{\angstrom^{-3}}$. With $\delta = \SI{0.3}{eV}$ and masses of a few tenths of $m_e$, say $0.5\,m_e$ for the holes and, per conduction valley, the density-of-states mass $(m^*_\parallel m^{*2}_\perp)^{1/3} = (0.92\times0.19^2)^{1/3}\,m_e = 0.32\,m_e$, the fractions are \num{5.3e-3} for one valence band and $6\times\num{2.7e-3} = \num{1.6e-2}$ for the six valleys. On the $40^3$ grid this leaves about 340 valence and 1{,}040 conduction points per band, so $3.5\times10^5$ pairs instead of $4.1\times10^9$, four orders of magnitude fewer. The window, not the interpolation, is what makes fine grids affordable near an edge.
\end{workedexample}

\subsection{Phonon linewidths and the Eliashberg function}

The phonon linewidth $\gamma^{\mathrm{ep}}_{\vect{Q}\nu}$ and the Eliashberg function $\alpha^2F(\omega)$ of \cref{sec:6.4} are Fermi-surface averages of $|g|^2$: double sums over $\vect{k}$ with two delta functions that place both the initial and the final electron on the Fermi surface. They are the hardest integrals in this section, because a delta function in three dimensions samples a two-dimensional surface, and a product of two samples the intersection of two surfaces. In practice each delta function is replaced by a Gaussian of width $\sigma$. The width must be larger than the spacing of the fine-grid energies near the Fermi surface, of order $|\nabla_{\vect{k}}\varepsilon|\,\Delta k$ with $\Delta k = |\vect{b}_i|/N_{\mathrm{fine}}$ the fine-grid step and $N_{\mathrm{fine}}$ the number of fine-grid points per direction, or the sum is noise. It must also be smaller than the energy scale on which $|g|^2$ and the density of states vary, or the result is biased. Convergence means decreasing $\sigma$ while refining the grid until the result is independent of both.

For metals with complicated Fermi surfaces the requirement is severe. For lead, the DFPT phonons needed $14^3$ $\vect{k}$ points and a $10^3$ wavevector grid with a Gaussian broadening of \SI{0.025}{Ry}, the electron--phonon vertices were interpolated from $6^3$ coarse grids, and the coupling strength $\lambda_{\mathrm{ep}}$ converged only with 400{,}000 quasi-random $\vect{k}$ points and 27{,}000 random wavevectors~\cite[Fig.~12]{ponce2016epw}. Quasi-random (Sobol) sampling converged faster than uniform grids in that study~\cite{ponce2016epw}, since it avoids the systematic aliasing of a regular grid with the Fermi surface. For an undoped insulator at $T = 0$ the electron--phonon linewidth below the gap vanishes (\cref{sec:6.4}), and the calculation is needed only for metals, superconductors and doped semiconductors.

\subsection{Convergence of the phonon-assisted dielectric function}

The phonon-assisted $\varepsilon_2^{\mathrm{PA}}(\omega)$ of \cref{sec:6.5} is a double sum over $\vect{k}$ and $\vect{Q}$ of the squared two-path amplitude, with one energy delta function~\cite[Eqs.~(1)--(3)]{noffsinger2012phonon}. Its convergence parameters are the two fine grids, the broadening of the delta function, the imaginary part $\varGamma_{\mathrm{w}}$ given to the intermediate states, and the number of intermediate bands. The energy resolution is set by the grid: in silicon, $40^3$ grids for both sums resolve features of \SI{14}{meV}, which is what the onset needs, while the smoother spectrum between the indirect and direct gaps converges at $24^3$ with a resolution of \SI{30}{meV}~\cite{noffsinger2012phonon}. In that range the result is insensitive to the intermediate-state width, set to a constant \SI{100}{meV}~\cite{noffsinger2012phonon}, because the intermediate states are far from resonance; near the direct gap they are not, and the width matters (\cref{sec:13.6}).

A coarse grid that looks converged for bands and phonons can still be too coarse for vertices, and a fine grid that resolves a \SI{30}{meV} spectrum can be too coarse for a \SI{14}{meV} onset. Each convergence statement belongs to the quantity, the energy range and the resolution for which it was made.

The comparison with experiment is where the physics outside the calculation shows. The computed spectra of Ref.~\cite{noffsinger2012phonon} had to be shifted rigidly by \SIrange{0.15}{0.23}{eV} to match the measured onsets. The shift absorbs the remaining error of the quasiparticle gap, of order \SI{0.1}{eV}, and the electron--phonon renormalization of the gap, \SIrange{0.05}{0.1}{eV} between 0 and \SI{400}{K}, which the calculation did not include~\cite{noffsinger2012phonon}. A computed absorption edge without such a shift is accurate only to the accuracy of the gap.

\subsection{A phonon-assisted electron-recoil calculation}

A phonon-assisted electron recoil (\cref{sec:13.2}) replaces the photon of \cref{sec:6.5} by the dark-matter density vertex at finite $\vect{q}$, and the four paths of its amplitude contain $\rho$, $g$ and the energies at three crystal momenta. The calculation chains the steps above, and each link has its own check. It starts from the ground state and DFPT on a coarse wavevector grid (\cref{app:E.4}), with the frequencies checked against neutron data and the Born charges against zero. Bands on a coarse $\vect{k}$ grid fine enough to place the conduction minimum fix the Wannier manifold, whose frozen window covers the final and intermediate states and whose bands are tested off the grid. The coarse-grid vertices $g$ are rotated into the Wannier gauge with their dipole and quadrupole parts removed, and tested against direct DFPT off the grid. For each target $\vect{q}$, $\rho_{mn}(\vect{k},\vect{q})$ is computed from the plane-wave coefficients and rotated with the same $U(\vect{k})$, and it too is tested against direct values off the grid. For each $\vect{q}$ the window is the range $\omega < E_{\mathrm{min}}(\vect{q})$ with no zero-phonon final state (\cref{sec:13.1}), cut to the fine-grid pockets of states that can contribute; its edges are in \cref{sec:13.9}. In that window the four-path amplitude is summed over the Wannier manifold, with the remainder above it from the projected Sternheimer equation \cref{eq:E.4-dpsiperp}, into the phonon-assisted loss function on a grid in $(|\vect{q}|,\omega)$.

Three controls close the chain. As $\vect{q}\to0$ the density vertex reduces to its optical limit (\cref{sec:3.2}), and the computed loss function must reproduce the phonon-assisted $\varepsilon_2^{\mathrm{PA}}$ of the same code and, after the gap shift above, the measured indirect absorption (\cref{sec:13.5}). With $g$ set to zero the phonon-assisted rate must vanish identically, which catches any zero-phonon contamination of the window. The third is the $f$-sum rule for the density vertices of one occupied state in its nonlocal form, \cref{eq:E.7-fsum}, whose shortfall measures the truncation of the intermediate-state sum once the nonlocal term $\delta_{\mathrm{NL}}$ is included.

The cost of a silicon run follows from the counts above. The DFPT stage is the one of \cref{we:E-dfpt-si}, a few hundred linear-response solutions. The density vertices need one non-self-consistent calculation per target $\vect{q}$; a grid of $|\vect{q}|$ in steps of $\sim\SI{0.1}{keV}$ up to a few keV and a few tens of directions in the irreducible wedge gives several hundred targets, and this stage dominates the plane-wave part of the work. The amplitude stage is a sum over pocket pairs. Near the edge the relevant states lie within $\sim\SI{0.1}{eV}$ of the band extrema, and the window fractions of \cref{we:E-counts}, which scale as $\delta^{3/2}$, drop by a factor $3^{3/2} \approx 5$ from their values at \SI{0.3}{eV}. On a $120^3$ fine grid this leaves about $1.8\times10^3$ valence points per band and $5.4\times10^3$ conduction points. Three valence bands lie within \SI{0.1}{eV} of the maximum, the heavy-hole, light-hole and split-off bands (\SI{44}{meV} below the others), so there are $\sim3\times10^7$ pairs per $\vect{q}$. Momentum conservation fixes the phonon wavevector of each pair, $\vect{Q} = \vect{k}' - \vect{k} - \vect{q}$ up to a reciprocal-lattice vector, so each pair needs its own interpolated vertices. Restricted to one external band and $N_W$ intermediate bands, the four paths need about $4\times216\times26\times6 \approx 1.3\times10^5$ complex multiply--adds of interpolation per pair, plus a smaller amount for $\rho$, which is $\sim\num{1e6}$ floating-point operations. That count assumes the $\vect{R}_p$ stage of \cref{eq:E.5-ginterp} is already done for the pair's $\vect{Q}$; it runs once per distinct $\vect{Q}$ of the pocket-difference set, at $216^2\times26\times6 \approx 7\times10^6$ multiply--adds per external band, and stays below the $\vect{R}_e$ stage for up to $\sim10^5$ distinct $\vect{Q}$. Per $\vect{q}$ the total is $\sim\num{3e13}$ operations, and for several hundred targets $\sim\num{1e16}$. At a sustained \num{e13} operations per second, typical of one modern accelerator, that is $10^3$\,s, a quarter of an hour; at \num{5e10}, typical of a few processor cores, it is two to three days. The estimate scales as $\delta^{3}$ with the energy range above the edge and as $N_{\mathrm{fine}}^6$ with the fine grid, so these two choices, not the plane-wave stages, decide the budget.
